\section{Exterior approximation and reduction to a rest end}
\label{n3:reduction:section}

We reduce the invariant-mass inequality to an energy inequality on a
controlled rest end.  We replace the distant end, bend the replacement
inside Schwarzschild spacetime, and repair its constraint deficit with a
conformal change of vanishing mass cost.  Throughout,
$(N,g,K)$ is a smooth exterior with one asymptotically flat end and compact
smooth boundary $B$.  The normal $\nu$ on $B$ points into $N$.  Completeness
means completeness up to this boundary.  We assume
\begin{equation}
\label{n3:reduction:original-decay}
 g-\delta=O_2(r^{-q}),\qquad K=O_1(r^{-1-q}),\qquad q>\tfrac12,
\end{equation}
and the integrability and dominant energy condition of
Theorem~\ref{n3:intro:exterior-inequality}.  The symbol $a_g(B)$ denotes the
filled enclosing perimeter infimum specified there.  None of the results
in this section requires the boundary to be connected or outermost.

\subsection{A strict conformal direction}

\begin{lemma}[Conformal constraint formulas]
\label{n3:reduction:conformal-formulas}
For a smooth function $f$, set
\[
 g_f=e^{4f}g,\qquad K_f=e^{2f}K.
\]
The corresponding constraint densities satisfy
\begin{align}
 16\pi e^{4f}\mu_f
   &=16\pi\mu-8\Delta_g f-8|df|_g^2,
       \label{n3:reduction:conformal-energy}\\
 8\pi J_f
   &=e^{-2f}\bigl(8\pi J+4K(\nabla f,\cdot)\bigr).
       \label{n3:reduction:conformal-momentum}
\end{align}
Here the second identity is an identity of covectors.  Consequently,
\begin{align}
 16\pi e^{4f}(\mu_f-|J_f|_{g_f})
 \ge{}&16\pi(\mu-|J|_g)\notag\\
 &+8\bigl(-\Delta_g f-|df|_g^2-|K|_g|df|_g\bigr).
       \label{n3:reduction:conformal-dec}
\end{align}
On the fixed boundary,
\begin{equation}
 \theta_{+,f}=e^{-2f}(\theta_++4\partial_\nu f).
       \label{n3:reduction:conformal-expansion}
\end{equation}
Equivalently, if $u=e^f>0$, the last summand in
Equation~\eqref{n3:reduction:conformal-dec} is
$8u^{-1}(-\Delta_g u-|K|_g|du|_g)$.
\end{lemma}

\begin{proof}
The scalar-curvature transformation in dimension three is
$e^{4f}\Scal_{g_f}=\Scal_g-8\Delta_g f-8|df|^2$.
Both quadratic terms in $K$ scale by $e^{-4f}$, which proves
Equation~\eqref{n3:reduction:conformal-energy}.
Put $\pi=K-(\tr_gK)g$.  Then
$\pi_f=e^{2f}\pi$.  For a covariant symmetric tensor $T$, the connection
difference for $\widetilde g=\Omega^2g$ gives, in dimension three,
\[
 \widetilde\nabla^{,i}(\Omega T)_{ij}
 =\Omega^{-1}\bigl(\nabla^iT_{ij}
          +2T_{ij}\nabla^i\log\Omega
          -(\tr_gT)\partial_j\log\Omega\bigr).
\]
Take $\Omega=e^{2f}$ and use $\tr_g\pi=-2\tr_gK$.
The terms containing $\tr_gK$ cancel, leaving
Equation~\eqref{n3:reduction:conformal-momentum}.  Taking its norm introduces
another factor $e^{-2f}$.  The triangle inequality, with the factor two
between the constraint normalizations $16\pi$ and $8\pi$, yields
Equation~\eqref{n3:reduction:conformal-dec}.
Finally, $\nu_f=e^{-2f}\nu$,
$H_f=e^{-2f}(H+4\partial_\nu f)$, and
$\tr_{B,g_f}K_f=e^{-2f}\tr_{B,g}K$.  This proves
Equation~\eqref{n3:reduction:conformal-expansion}.  The last assertion follows
from $\Delta\log u+|d\log u|^2=u^{-1}\Delta u$.
\end{proof}

\begin{lemma}[Strict direction under the original decay]
\label{n3:reduction:strict-direction}
Choose
\[
 0<\delta<\beta<\min(q,1).
\]
There are smooth positive functions $w$ and $\phi$ on $N$, with
$w=r^{-3-\delta}$ on the end, such that
\begin{align}
 &\partial_\nu\phi=-1\quad\hbox{on }B,
       &\phi&=O_2(r^{-1}),\label{n3:reduction:strict-boundary}\\
 &-\Delta_g\phi-|K|_g|d\phi|_g\ge w,
       &\frac{-\Delta_g\phi}{w}&\longrightarrow1.
       \label{n3:reduction:strict-source}
\end{align}
The Euclidean normal flux
\[
 L_\phi=\lim_{r\to\infty}
          \int_{S_r}\partial_r\phi\,\dd A_\delta
\]
is finite and negative.  For every finite $s\ge0$, set
\begin{equation}
 \label{n3:reduction:linear-family}
 u_s=1+s\phi,\qquad g_s=u_s^4g,\qquad K_s=u_s^2K.
\end{equation}
These data are complete up to $B$ and satisfy DEC, strictly for $s>0$.
If $\theta_+\le0$ on $B$, their boundary expansion remains nonpositive
and is strictly negative for $s>0$.  Their decay exponent can be taken to be
$\min(q,1)>1/2$, their constraint densities are integrable, and
\begin{equation}
 \label{n3:reduction:strict-charges}
 E_s=E+sA_\phi,\qquad P_s=P,
 \qquad A_\phi=-\frac{L_\phi}{2\pi}>0.
\end{equation}
Their full enclosing infima satisfy
\begin{equation}
 \label{n3:reduction:strict-areas}
 a_g(B)\le a_{g_s}(B)
       \le (1+s\|\phi\|_\infty)^4a_g(B).
\end{equation}
In particular $E_s\to E$ and $a_{g_s}(B)\to a_g(B)$ as $s\downarrow0$.
No sign assumption on the ADM vector is needed.

If $K$ is compactly supported and the metric has symbol estimates of all
orders with $g-\delta=O(r^{-1})$, the function can instead be chosen with
$\phi=O_k(r^{-1})$ for every $k$, while retaining
Equations~\eqref{n3:reduction:strict-boundary} and
\eqref{n3:reduction:strict-source}.
\end{lemma}

\begin{proof}
Choose smooth $b\ge|K|_g$, positive everywhere and equal to
$b_0r^{-1-\beta}$ sufficiently far out, where $b_0>0$ is chosen large
enough.  This is possible because $\beta<q$.  Choose a smooth positive
regularizer $\zeta$, equal to $r^{-2}$ on the end, and extend $w$ smoothly
and positively to $N$.  We solve
\begin{equation}
 \label{n3:reduction:drift-poisson}
 -\Delta_g\phi=b\sqrt{|d\phi|_g^2+\zeta^2}+w,
 \qquad \partial_\nu\phi=-1,
 \qquad \phi\longrightarrow0.
\end{equation}
The regularizer makes this equation smooth even at critical points.

First truncate at a large coordinate sphere $S_R$ and prescribe
$\phi=0$ there.  The Neumann and Dirichlet conditions occur on disjoint
boundary components.  For $0\le t\le1$, replace $b$ in
Equation~\eqref{n3:reduction:drift-poisson} by $tb$.  At $t=0$, the mixed
Laplacian is invertible: its Dirichlet face gives the Poincare inequality,
and the weak solution obtained from its coercive quadratic form is smooth
by boundary regularity.  At any solution, the linearization is a
Laplacian with a smooth drift of norm at most $b$ and with the homogeneous
mixed boundary conditions.  The strong maximum principle and the Hopf
boundary principle make its kernel zero.  It is a Fredholm perturbation
of the mixed Laplacian of index zero, and is therefore invertible.
The boundary estimates used here are the ordinary elliptic estimates
for Dirichlet and normal Neumann conditions: positive definiteness of
$g$ verifies their principal complementing condition, and the disjoint
boundary components can be covered separately
\cite{AgmonDouglisNirenberg1959}.  The maximum principles and local
elliptic estimates below are used in their usual uniformly elliptic
forms \cite{GilbargTrudinger2001}.

To close continuation at $t=1$, first work on a fixed truncation.
The elliptic $W^{2,p}$ estimate
and gradient interpolation give
\[
 \|\phi\|_{W^{2,p}}
 \le C_p\bigl(1+\|\phi\|_{L^p}+\|d\phi\|_{L^p}\bigr)
 \le \tfrac12\|\phi\|_{W^{2,p}}
       +C'_p(1+\|\phi\|_{L^p}).
\]
The constants include the fixed boundary data and are uniform in $t$.
If the supremum norms of solutions on this truncation were unbounded,
divide them by their supremum norms.  The preceding estimate with
$p>3$, compactness, and the equation give a nonzero limit $v$ with
homogeneous mixed data satisfying
\[
 -\Delta_gv=t_*b|dv|_g.
\]
This can be written $\Delta_gv+A\cdot dv=0$ with a bounded measurable
drift: take $A=t_*b\nabla v/|dv|$ off the critical set and zero on it.
The strong and Hopf maximum principles exclude a nonzero solution with
the homogeneous mixed conditions.  Thus the supremum norms are bounded.
The $W^{2,p}$ estimate, Schauder estimates, and the smooth equation now
give the bounds needed for closedness of continuation.  Existence on
each truncation follows.

Every such solution is nonnegative.  A negative minimum cannot be
interior because the right side of
Equation~\eqref{n3:reduction:drift-poisson} is positive.  At an inner-boundary
minimum, the outward-domain derivative would be negative by the Hopf
principle, whereas the prescribed derivative in that direction is $1$.
The outer value is zero.

It remains to make these bounds independent of $R$.  For fixed sufficiently large
$A$ and then sufficiently large $r_0$, the positive radial function
\[
 v_0(r)=r^{-1}(1-Ar^{-\delta})
\]
satisfies, on $r\ge r_0$,
\begin{align*}
 -\Delta_gv_0-b|dv_0|_g
 &=A\delta(1+\delta)r^{-3-\delta}
       +O(r^{-3-q})+O(r^{-3-\beta})\ge c_0w.
\end{align*}
The constants are fixed independently of the truncation.  Since
$b\zeta=o(w)$, a sufficiently large multiple of $v_0$ is a supersolution
of the regularized equation.  The comparison principle gives
\begin{equation}
 \label{n3:reduction:core-barrier}
 0\le\phi(x)\le C(1+M_R)v_0(r),\qquad
 M_R=\sup_{N\cap\{r\le r_0\}}\phi,
 \quad r_0\le r\le R.
\end{equation}
The notation for the compact core includes all points outside the end
chart.  Were $M_R$ unbounded along expanding truncations, the normalized
solutions $\phi/M_R$ would have locally uniform $W^{2,p}$ bounds, including
at $B$, by the preceding local estimates and
Equation~\eqref{n3:reduction:core-barrier}.  A subsequence would converge
locally in $C^1$ to a nonnegative function with maximum one on the core,
homogeneous inner Neumann data, and a bounded-drift homogeneous equation.
Equation~\eqref{n3:reduction:core-barrier} makes its value tend to zero at
infinity.  It consequently attains a positive maximum, contradicting the
strong or Hopf principle for that homogeneous equation.  This proves the
uniform core bound.  The contradiction uses the homogeneous bounded-drift
equation.

Local compactness and diagonal extraction give a smooth solution of
Equation~\eqref{n3:reduction:drift-poisson}, with $\phi=O(r^{-1})$.
The claimed derivative count can be checked directly on annuli.  Put
\[
 \phi_{\rho}(y)=\rho\phi(\rho y),\qquad
 g_{\rho}(y)=g(\rho y),\qquad 1<|y|<4.
\]
On smaller fixed annuli the equation is
\[
 -\Delta_{g_{\rho}}\phi_{\rho}
 =\rho b(\rho y)
       \sqrt{|d\phi_{\rho}|_{g_{\rho}}^2+
                     (\rho^2\zeta(\rho y))^2}
       +\rho^3w(\rho y).
\]
Its coefficients have uniform $C^{1,1}$ bounds from the two derivatives
in Equation~\eqref{n3:reduction:original-decay}; the scaled drift is
$O(\rho^{-\beta})$ and the scaled source is $O(\rho^{-\delta})$.
The bounded zeroth-order norm, interior $W^{2,p}$ estimates, and gradient
interpolation first give uniform $C^{1,\alpha}$ bounds for some
$\alpha>0$.  The right side is then uniformly $C^\alpha$.  Schauder
estimates give uniform $C^{2,\alpha}$ bounds on smaller annuli.  This
proves $\phi=O_2(r^{-1})$ without using a third derivative of $g$ or a
second derivative of $K$.

The right side of Equation~\eqref{n3:reduction:drift-poisson} is
$w+O(r^{-3-\beta})$.  It is integrable and proves
Equation~\eqref{n3:reduction:strict-source}.  The divergence theorem, with
outward-domain normal $-\nu$ on $B$, gives
\begin{equation}
 \label{n3:reduction:phi-flux}
 \lim_{r\to\infty}\int_{S_r}\partial_{n_g}\phi\,\dd A_g
 =-|B|_g-\int_N
       \bigl(b\sqrt{|d\phi|_g^2+\zeta^2}+w\bigr)\,\dd V_g.
\end{equation}
The Euclidean flux differs by $O(r^{-q})$ and has the same finite limit.

For every fixed finite $s\ge0$, the factor $u_s$ in
Equation~\eqref{n3:reduction:linear-family} is bounded and at least one,
so the transformed data are complete up to $B$.
The $u$-form of Lemma~\ref{n3:reduction:conformal-formulas} gives
\begin{align*}
 16\pi u_s^4(\mu_s-|J_s|_{g_s})
 &\ge16\pi(\mu-|J|_g)
        +\frac{8s}{u_s}(-\Delta_g\phi-|K|_g|d\phi|_g)
 \ge \frac{8s}{u_s}w,\\
 \theta_{+,s}
 &=u_s^{-2}\left(\theta_+-\frac{4s}{u_s}\right).
\end{align*}
This proves the asserted signs for all finite $s$, without a smallness
restriction.

Since $u_s-1=O_2(r^{-1})$, the transformed decay exponent is
$\min(q,1)>1/2$.  The exact energy transformation is
\[
 16\pi u_s^4\mu_s=16\pi\mu-8s u_s^{-1}\Delta_g\phi.
\]
Here $\Delta_g\phi=O(r^{-3-\delta})$, and the additional current term
is bounded by a fixed-$s$ constant times
$|K|_g|d\phi|_g=O(r^{-3-q})$.  Both are integrable in dimension three.
The bounded positive conformal factor also preserves integrability in
the transformed volume measure.

The leading metric change is $4s\phi\delta$, whose ADM energy flux
is $-8s\int_{S_r}\partial_r\phi\,\dd A_\delta$ before division by
$16\pi$.  All remaining metric flux terms are
$O_s(r^{-q})+O_s(r^{-1})$.  Writing $\pi=K-(\tr_gK)g$, one has
$\pi_s=u_s^2\pi$, so the change in its Euclidean momentum flux is
$O_s(r^2r^{-1}r^{-1-q})=O_s(r^{-q})$.
This proves Equation~\eqref{n3:reduction:strict-charges};
Equation~\eqref{n3:reduction:phi-flux} gives $A_\phi>0$.
No $1/r$ coefficient expansion for $\phi$ is needed.
Finally, the cut class is unchanged and every cut has area
$\int_\Gamma u_s^4\,\dd A_g$.  Bounding $u_s$ between $1$ and
$1+s\|\phi\|_\infty$ and taking infima proves
Equation~\eqref{n3:reduction:strict-areas}, including coincident and
disconnected cuts.

For the final assertion, take $b$ nonnegative, compactly supported, and
at least $|K|$; the proof is unchanged.  The end equation is then simply
$-\Delta_g\phi=w$.  Once the two-derivative estimate has been established,
successive annular elliptic estimates use the assumed symbol bounds for
$g$ and $w$ to give $\phi=O_k(r^{-1})$ for every $k$.
\end{proof}

\subsection{Compactness of enclosing minimizers}

The end replacement will change the metric on receding annuli.  The next
lemma keeps area-minimizing enclosures in a fixed compact set, where local
metric convergence controls their areas.

\begin{lemma}[Enclosing areas under receding changes]
\label{n3:reduction:area-compactness}
Let $g_j$ be smooth metrics on the fixed exterior $N$, converging to $g$
locally uniformly up to $B$.  Suppose that in a fixed chart outside a
compact set they satisfy, with constants independent of $j$,
\begin{equation}
 \label{n3:reduction:area-geometry}
 c\delta\le g_j\le C\delta,\qquad r|\partial g_j|_\delta\le C.
\end{equation}
Suppose also that each $g_j$ has an asymptotically flat end, possibly
in a different chart, and admits a fixed compact enclosing competitor
with uniformly bounded area.  Then minimizing filled enclosures exist,
their boundaries lie in one fixed compact subset of $N$, and
\[
 a_{g_j}(B)\longrightarrow a_g(B).
\]
The perimeter and smooth enclosing infima agree.  Each minimizing
boundary is $C^{1,1}$, and is smooth minimal away from contact with $B$.
\end{lemma}

\begin{proof}
For purposes of perimeter minimization only, extend the metrics through
a collar of $B$ and attach a fixed compact filling.  The filling need not
satisfy any constraint or curvature condition.  Extend $g_j$ so that the
extensions converge uniformly on this fixed collar and filling.  A filled
enclosure is then a bounded set containing the prescribed inner obstacle.

For each fixed metric, sufficiently large spheres in its asymptotically
flat chart have strictly positive outward mean curvature.  Clipping a
bounded enclosure at any such sphere cannot increase perimeter.  To see
this, integrate the divergence of the outward unit normal field of the
sphere foliation over the part of the enclosure outside the sphere.  Its
divergence is positive.  Its flux through the exterior boundary is at
most the area of that boundary, whereas its flux on the cutting sphere
is the negative of the new cutting area.  Thus the discarded boundary
area is at least the inserted area.  Approximation gives the same
statement for finite-perimeter sets.  The direct method on the resulting
bounded domain now gives a minimizer.  The clipping observation makes
this a minimizer against every bounded competitor, not merely those
inside the chosen truncation.

Off the obstacle the boundary is locally perimeter minimizing.  If it
contains a point of radius $r$ sufficiently large, the coordinate ball
of radius $c_1r$ about that point avoids the obstacle.  After scaling
this ball to unit size, Equation~\eqref{n3:reduction:area-geometry} gives
uniform ellipticity and uniform first-derivative bounds for its area
integrand.  The local perimeter density estimate gives
\[
 \operatorname{Area}_{g_j}
       (\partial F_j\cap B_{c_1r})\ge c_2r^2.
\]
Here $c_1,c_2>0$ are independent of $j$ and $r$.  One may obtain this
estimate from the relative isoperimetric inequality and the comparison
which removes, or fills, the set in concentric balls; integration of the
resulting differential inequality gives the area lower bound.  The
fixed competitor bounds the left side above independently of $j$.
Consequently the minimizing boundaries lie in a fixed compact set.

Perimeter compactness now gives a limiting filled enclosure.  Uniform
metric convergence on that compact set and lower semicontinuity imply
\[
 a_g(B)\le\liminf_j a_{g_j}(B).
\]
Conversely, using a fixed enclosure with $g$-area arbitrarily close to
$a_g(B)$ gives the reverse upper-limit inequality.

The local obstacle regularity theorem for pure area minimization with
a $C^2$ obstacle gives $C^{1,1}$ regularity, and smoothness off contact,
in dimension three; its local estimates require only the appropriate
bounded geometry of the smooth metric
\cite[Regularity Theorem 1.3]{HuiskenIlmanen2001}.
This applies to the fixed smooth obstacle and our smooth one-sided
metrics after the collar extension.  Finally, push a minimizing boundary
slightly away from its contact set by a smooth vector field agreeing
there with the outward obstacle normal.  In a small obstacle collar its
flow increases the signed distance to the obstacle.  The resulting
$C^1$ embedded boundary is a positive distance from the obstacle.
Smooth local graph approximation, within that distance, preserves
enclosure and changes area by an arbitrarily small amount.  This proves
equality of the smooth and perimeter infima and also justifies using
smooth fixed competitors in the upper-limit argument.
\end{proof}

\subsection{Schwarzschild reference charges}

We next construct the vacuum end used in the replacement.  Its charges
must agree with those of the original data before we bend it to a rest
slice.

\begin{lemma}[A reference end with prescribed timelike charges]
\label{n3:reduction:boosted-reference}
Let $E>|P|$ and $m=(E^2-|P|^2)^{1/2}$.  There is a spacelike plane
near spatial infinity of Schwarzschild spacetime of mass $m$ whose
induced data, in asymptotically Euclidean coordinates on the plane,
have ADM charges $(E,P)$.  The reference data satisfy
\[
 g_*-\delta=O_k(r^{-1}),\qquad K_*=O_k(r^{-2})
\]
for every derivative order $k$.
\end{lemma}

\begin{proof}
In static isotropic coordinates $(T,y)$ the Schwarzschild metric is
\begin{equation}
 \label{n3:reduction:schwarzschild-spacetime}
 -\left(\frac{1-m/(2r)}{1+m/(2r)}\right)^2\dd T^2
       +\left(1+\frac{m}{2r}\right)^4\delta_{ij}\dd y^i\dd y^j,
 \qquad r=|y|.
\end{equation}
Its perturbation $h$ from Minkowski space has symbol decay $r^{-1}$
on every fixed spacelike cone.  Constant Lorentz transformations preserve
this decay.  We verify the transformation of the ADM charges directly.

In inertial coordinates for the background Minkowski metric $\eta$, put
\begin{align}
 \mathcal F_{\ell ab}
  ={}&\partial_a h_{\ell b}-\partial_\ell h_{ab}\notag\\
   &+\eta_{\ell b}
       (\partial^d h_{ad}-\partial_a\tr_\eta h)
    -\eta_{ab}
       (\partial^d h_{\ell d}-\partial_\ell\tr_\eta h).
       \label{n3:reduction:superpotential}
\end{align}
This tensor is antisymmetric in its first two slots, and differentiation
of the displayed formula gives
$\partial^\ell\mathcal F_{\ell ab}=2G^{\mathrm{lin}}_{ab}$.
The exact Einstein tensor vanishes; its terms beyond the linearization
are $O(r^{-4})$.  Thus this divergence is $O(r^{-4})$.

Use an inertial frame in which the plane under consideration is
$x^0=0$.  With $\eta=\operatorname{diag}(-1,1,1,1)$, direct substitution
in Equation~\eqref{n3:reduction:superpotential} gives
\begin{align*}
 \mathcal F_{i00}&=\partial_jh_{ij}-\partial_ih_{jj},\\
 \mathcal F_{i0k}
  &=2\bigl(K^{\mathrm{lin}}_{ik}
          -\delta_{ik}\tr K^{\mathrm{lin}}\bigr)
       +\partial_kh_{0i}-\delta_{ik}\partial_jh_{0j},
\end{align*}
where
\[
 K^{\mathrm{lin}}_{ik}
 =\tfrac12(\partial_0h_{ik}-\partial_ih_{0k}-\partial_kh_{0i}).
\]
The convention is the normal rate of the spatial metric divided by two.
For fixed $k$, the final vector has zero flux through every closed sphere.
Indeed it equals
\[
 \partial_j(\delta_{jk}h_{0i}-\delta_{ik}h_{0j}),
\]
the divergence of a tensor antisymmetric in $i,j$.  This zero-flux
identity is local to a neighborhood of the sphere: extend the potential
smoothly inside before applying the divergence theorem, if necessary.
The nonlinear errors in the induced metric and second fundamental form
give vanishing flux errors.  Therefore the limiting fluxes of
$\mathcal F_{i0b}$ are $16\pi(E,P)$.

These fluxes are the integrals of the dual two-form of the first two
slots of $\mathcal F$, with its translation slot held fixed.  They
therefore transform tensorially under a constant Lorentz change of
frame.  Independence of the limiting spacelike cut follows directly
from Stokes' theorem: two large cuts on planes of slopes bounded
strictly below one can be joined within a fixed spacelike cone by a
three-dimensional cylinder of volume $O(R^3)$, all of whose points have
radius comparable to $R$.  Its divergence error is $O(R^{-4})$, so the
flux difference is $O(R^{-1})$.  Replacing the resulting ellipsoidal cuts
by coordinate spheres on either plane gives the same estimate.  The
limiting translation charges consequently transform by the Lorentz
representation.  The static plane has charges $(m,0)$, as follows by
inserting Equation~\eqref{n3:reduction:schwarzschild-spacetime} into the ADM
integral.  Its Lorentz orbit is precisely the set of future timelike
vectors of norm $m$, proving the assertion.

More explicitly the plane can be written $T=v\cdot y$, $|v|<1$.
Writing $y=Ax$, where
$A=(I-v\otimes v)^{-1/2}$, makes the induced Minkowski metric in the
$x$ coordinates equal to $\delta$.  The displayed decay then follows
from the symbol estimates in
Equation~\eqref{n3:reduction:schwarzschild-spacetime} and the formula for
the second fundamental form.  These coordinates will be used when
identifying the reference end with the given end.
\end{proof}

\subsection{Moments and compact linear corrections}

All operators in the next two lemmas are Euclidean.  For symmetric
covariant tensors define
\begin{equation}
 \label{n3:reduction:linear-operators}
 \mathcal Lh=\partial_i\partial_jh_{ij}-\Delta\tr h,
 \qquad
 (\mathcal Dk)_i=\partial_j(k_{ij}-(\tr k)\delta_{ij}).
\end{equation}
The formal adjoint kernel of $\mathcal L$ consists of affine functions;
the formal adjoint kernel of $\mathcal D$ consists of Euclidean Killing
fields.  Orthogonality to these kernels is the compatibility condition
for the compact inverses below.

\begin{lemma}[Compact scalar and symmetric divergence inverses]
\label{n3:reduction:compact-inverses}
Let $\mathcal A\subset\R^3$ be a fixed connected open annulus and let
$Q\Subset\mathcal A$.  Smooth sources supported in $Q$ admit the
following compact corrections, supported in a fixed larger compact
subset of $\mathcal A$.
\begin{enumerate}[label=\textup{(\roman*)}]
\item If $\int f=0$ and $\int x^if=0$ for $i=1,2,3$, there is a smooth
      symmetric tensor $h$ with $\mathcal Lh=f$ and
      $\|h\|_{W^{k+2,p}}\le C_{k,p}\|f\|_{W^{k,p}}$.
\item If $\int V\cdot Y=0$ for every Euclidean Killing field $Y$, there
      is a smooth symmetric tensor $k$ with $\mathcal Dk=V$ and
      $\|k\|_{W^{k_0+1,p}}\le C_{k_0,p}\|V\|_{W^{k_0,p}}$.
\end{enumerate}
The estimates hold for integers $k,k_0\ge0$ and $1<p<\infty$.
\end{lemma}

\begin{proof}
We first record the ordinary divergence inverse being used.  On a ball,
the regularized Bogovskii operator solves $\Div u=f$ for a mean-zero
compactly supported source, preserves compact support, and gains one
derivative in $W^{k,p}$; see
\cite[\S3.1, Equation~(3.13), Theorem~3.2, Corollary~3.4, and Remark~3.5]{CostabelMcIntosh2010}.
To apply it on the annulus, cover $Q$ by finitely
many balls compactly contained in $\mathcal A$, adding finitely many
overlapping balls so that the overlap graph is connected.  Partition the
source among the balls.  On a spanning tree of the overlap graph, pass
each piece's integral to its parent by subtracting and adding that
integral times a fixed unit-integral bump in the overlap.  Starting at
the leaves leaves every piece with zero integral; the root does also
because the total integral vanishes.  The decomposition has bounded
Sobolev norms because all cutoffs and bumps are fixed.  The sum of the
ball inverses gives a compact inverse on the annulus.  For sources in a
fixed compact set its output support lies in another fixed compact set.
This construction may be repeated twice using slightly enlarged compact
sets still contained in $\mathcal A$.

For the scalar assertion, solve $\partial_i u_i=f$ compactly.
Integration by parts gives
\[
 \int u_i=-\int x^if=0.
\]
Apply the divergence inverse to each component of $u$, obtaining
$\partial_j A_{ij}=u_i$.  Symmetrizing $A$ does not change its double
divergence.  Denote this symmetric matrix by $S$ and put
\[
 h=S-\tfrac12(\tr S)\delta.
\]
Since $h-(\tr h)\delta=S$, its linearized scalar curvature is $f$.
The two applications of the divergence inverse give the two-derivative
gain.

For the vector assertion, translation orthogonality gives
$\int V_i=0$.  Solve $\partial_jB_{ij}=V_i$ componentwise.  Orthogonality
to rotations and integration by parts imply
\[
 0=\int(x^jV_i-x^iV_j)
   =-\int(B_{ij}-B_{ji}).
\]
Thus the skew part $S_{ij}=(B_{ij}-B_{ji})/2$ has zero integral.  Apply
the divergence inverse to its independent components to obtain a tensor
$D$, skew in its first two indices, such that
$\partial_kD_{ijk}=-S_{ij}$.  Define
\begin{equation}
 \label{n3:reduction:stress-symmetrization}
 C_{ijk}=D_{ijk}-D_{ikj}-D_{jki}.
\end{equation}
Index interchange shows both
$C_{ijk}=-C_{ikj}$ and
$(C_{ijk}-C_{jik})/2=D_{ijk}$.  Consequently
\[
 T_{ij}=B_{ij}+\partial_kC_{ijk}
\]
is symmetric and satisfies $\partial_jT_{ij}=V_i$.  The last statement
uses the skew symmetry in $j,k$ to cancel the double divergence of $C$.
Set $k=T-(\tr T)\delta/2$ to invert trace reversal.  The first inverse
gains one derivative, the second gains another, and the last derivative
of $C$ leaves a net gain of one.  All constructions preserve smoothness
and have the stated support and norm bounds.
\end{proof}

\begin{lemma}[Integrable sources give sublinear first moments]
\label{n3:reduction:moment-estimates}
For data satisfying Equation~\eqref{n3:reduction:original-decay} and
integrable $\mu,|J|$, put $h=g-\delta$ and
$\pi=K-(\tr_\delta K)\delta$.  Then
\begin{align}
 \mathcal Lh&=16\pi\mu+O(r^{-2-2q}),\notag\\
 \partial_j\pi_{ij}&=8\pi J_i+O(r^{-2-2q}).
       \label{n3:reduction:linear-source-integrability}
\end{align}
In particular these Euclidean sources are integrable.  Their scalar
boundary fluxes paired with affine functions having zero constant term,
and their vector boundary fluxes paired with rotational Killing fields,
are $o(R)$ on coordinate spheres of radius $R$.
\end{lemma}

\begin{proof}
The scalar-curvature expansion has remainder bounded by
\[
 C\bigl(|h||\partial^2h|+|\partial h|^2+|K|^2\bigr),
\]
and the replacement of the momentum constraint by Euclidean divergence
has remainder bounded by
$C(|h||\partial K|+|\partial h||K|)$.  These are
$O(r^{-2-2q})$, using exactly the derivatives in
Equation~\eqref{n3:reduction:original-decay}.  They are integrable because
$2+2q>3$.

For any integrable function $F$ on the end and fixed $c>1$,
\begin{equation}
 \label{n3:reduction:sublinear-tail}
 \frac1R\int_{r_0<|x|<cR}|x|\,|F(x)|\,\dd x\longrightarrow0.
\end{equation}
Indeed the contribution inside any fixed radius $M$ tends to zero
after division by $R$; the remaining contribution is at most
$c\int_{|x|>M}|F|$.  Let $R\to\infty$ and then $M\to\infty$.

For an affine function $f$, the vector field
\[
 \mathcal Q_f(h)^i
  =f(\partial_jh_{ij}-\partial_i\tr h)
       -(\partial_jf)h_{ij}+(\partial_if)\tr h
\]
satisfies $\partial_i\mathcal Q_f(h)^i=f\mathcal Lh$.
For a Euclidean Killing field $Y$,
\[
 \partial_j(\pi_{ij}Y^i)=Y^i\partial_j\pi_{ij},
\]
because $\pi$ is symmetric.  Integrating from a fixed inner sphere and
using Equation~\eqref{n3:reduction:sublinear-tail} proves the asserted
$o(R)$ estimates.  This argument does not assert the existence of a
center-of-mass or angular-momentum limit and requires no parity
condition.
\end{proof}

\begin{proposition}[Annular replacement with a vanishing deficit]
\label{n3:reduction:annular-replacement}
Suppose $E>|P|$, and let $(g_*,K_*)$ be the reference end from
Lemma~\ref{n3:reduction:boosted-reference}, identified with the given end
in the common asymptotically Euclidean coordinates.  For all sufficiently
large $R$ there are smooth data $(\widetilde g_R,\widetilde K_R)$ equal to
$(g,K)$ on $r\le R$ and to $(g_*,K_*)$ on $r\ge4R$ such that
\begin{equation}
 \label{n3:reduction:annular-deficit}
 16\pi(\widetilde\mu_R-|\widetilde J_R|_{\widetilde g_R})
       \ge-\eta_RR^{-3},\qquad \eta_R\longrightarrow0.
\end{equation}
The possible negative part is supported in $R<r<4R$.
The metrics there are uniformly Euclidean comparable and have uniformly
bounded scaled first derivatives.  The construction uses no derivatives
of the original data beyond those in
Equation~\eqref{n3:reduction:original-decay} for its quantitative estimates.
\end{proposition}

\begin{proof}
Fix $\tfrac12<\beta<\min(q,1)$ and work on the annulus
$\mathcal A=\{1<|y|<4\}$ with $x=Ry$.  The scaled fields are
\[
 g_R(y)=g(Ry),\qquad k_R(y)=RK(Ry),
\]
and likewise for the reference data.  Their constraint densities are the
original densities multiplied by $R^2$.  If $h_R=g_R-\delta$, then
\begin{equation}
 \label{n3:reduction:scaled-smallness}
 \|h_R\|_{C^2(\mathcal A)}+\|k_R\|_{C^1(\mathcal A)}
 +\|h_{*,R}\|_{C^2(\mathcal A)}+\|k_{*,R}\|_{C^1(\mathcal A)}
 \le CR^{-\beta}.
\end{equation}
Choose $\chi\in C^\infty(\mathcal A)$ with $0\le\chi\le1$, equal to one near $|y|=1$ and zero
near $|y|=4$, with all its derivatives supported in a fixed compact
subannulus.  Blend both fields by this cutoff.

Write the constraints in the unnormalized form
\[
 \mathcal C(g,k)
 =\bigl(\Scal_g+(\tr_gk)^2-|k|_g^2,
             \Div_g(k-(\tr_gk)g)\bigr).
\]
Their Euclidean linear part is $(\mathcal Lh,\mathcal Dk)$.
On fields satisfying Equation~\eqref{n3:reduction:scaled-smallness}, the
nonlinear remainder is bounded in supremum norm by $CR^{-2\beta}$.
Thus the blended constraints equal the blended original constraint
sources, plus a quadratic error and the compact commutators
\begin{align*}
 f_R&=\mathcal L(\chi(h_R-h_{*,R}))
                      -\chi\mathcal L(h_R-h_{*,R}),\\
 V_R&=\mathcal D(\chi(k_R-k_{*,R}))
                      -\chi\mathcal D(k_R-k_{*,R}).
\end{align*}
They have $C^{0,1}$ norms at most $CR^{-\beta}$.  In detail the scalar
commutator applied to a symmetric tensor $h$ is
\begin{align*}
 &2\partial_i\chi(\partial_jh_{ij}-\partial_i\tr h)
       +(\partial_i\partial_j\chi)h_{ij}-(\Delta\chi)\tr h;
\end{align*}
one derivative of this expression uses at most two derivatives of $h$.
The vector commutator is $(\partial_j\chi)(k_{ij}-(\tr k)\delta_{ij})$;
one derivative uses at most one derivative of $k$.  No derivative of the
complete matter sources is needed.

Every affine moment of $f_R$ and every Killing-field moment of $V_R$
is $o(R^{-1})$.  For the constant moments, integration by parts on
$\mathcal A$ gives $R^{-1}$ times the difference between the physical
translation fluxes on its two boundary spheres, less the integrals of
the blended linear sources.  The flux difference tends to zero because
the reference and original ADM vectors agree.  Replacing the physical
momentum integrand by its Euclidean trace reversal costs
$O(R^{1-2q})$, which tends to zero.  By
Lemma~\ref{n3:reduction:moment-estimates}, the physical linear-source
integrals on this receding annulus tend to zero as well.
For a linear scalar test or a rotational vector test the corresponding
factor is $R^{-2}$.  Its physical boundary fluxes are $o(R)$ by the same
lemma, while the weighted source integrals are $o(R)$ by
Equation~\eqref{n3:reduction:sublinear-tail}.  This proves the claim.

Remove these small moments by fixed smooth bumps.  Explicitly, for a
basis $p_1,\ldots,p_4$ of affine functions and a fixed nonnegative bump
$\omega$ positive on a ball inside the annulus, the Gram matrix
\[
 G_{ab}=\int\omega p_ap_b
\]
is positive definite.  Subtract
$\omega\sum_{a,b}p_a(G^{-1})_{ab}\int p_bf_R$ from $f_R$.
For the vector moments use in exactly the same way the Gram matrix
$\int\omega Y_a\cdot Y_b$ of a basis of the six Euclidean Killing fields.
It too is positive definite, since a nonzero Killing field cannot vanish
on an open ball.  Both bump corrections have every fixed smooth norm
$o(R^{-1})$.

Apply Lemma~\ref{n3:reduction:compact-inverses} with the negative of the
moment-free commutators.  Since their $W^{1,p}$ norms are $O(R^{-\beta})$,
choosing $p>3$ gives metric and tensor corrections of size $O(R^{-\beta})$
in $W^{3,p}$ and $W^{2,p}$, respectively.  Sobolev embedding gives the
needed $C^{2,\alpha}$ and $C^{1,\alpha}$ bounds.  Their supports stay in
a fixed compact subset of $\mathcal A$.  Extend them by zero and undo
the scaling.  The metric remains positive definite for large $R$.

To check DEC, denote the original scaled sources by $(A_R,B_R)$, so
$A_R\ge2|B_R|_{g_R}$.  The reference sources vanish.  The corrected
sources therefore have the form
\[
 \bigl(\chi A_R,\chi B_R\bigr)
       +o(R^{-1})+O(R^{-2\beta})
\]
in supremum norm.  The norm of $B_R$ changes by at most
$CR^{-\beta}|B_R|\le CR^{-2\beta}$ when the metric is changed, since
the pointwise constraint formula gives $|B_R|\le CR^{-\beta}$.
All other nonlinear correction terms have the same quadratic bound.
Because $2\beta>1$, the scaled DEC deficit is $o(R^{-1})$.
Multiplication by $R^{-2}$ gives
Equation~\eqref{n3:reduction:annular-deficit}.  Outside the correction annulus
the data are either the original DEC data or the exact vacuum reference
data.  Smoothness follows from smoothness of the original fields and of
the compact inverse constructions, irrespective of the sizes of their
unneeded higher derivatives.
\end{proof}

\subsection{Bending to a rest slice and repairing the deficit}

\begin{lemma}[A vacuum bend with uniform geometry]
\label{n3:reduction:vacuum-bend}
The data in Proposition~\ref{n3:reduction:annular-replacement} can be
modified farther out, inside their exact Schwarzschild part, so that
they are induced by a constant-static-time slice outside a compact set.
This modification adds no constraint deficit.  Write $(g_R,K_R)$ for
the resulting unscaled data.  In the common end
coordinates the metrics are uniformly Euclidean comparable,
and throughout all transition annuli
\begin{equation}
 \label{n3:reduction:bend-geometry}
 |\partial g_R|_\delta+|K_R|_\delta\le C/r.
\end{equation}
There is a smooth proper radius $\rho_R$ on the end which equals the original
coordinate radius on the gluing annulus and equals the rest isotropic
radius beyond the transitions, with
\begin{equation}
 \label{n3:reduction:radius-geometry}
 c r\le\rho_R\le C r,\qquad
 |d\rho_R|_{g_R}^2\ge c,
 \qquad |\Delta_{g_R}\rho_R|\le C/r.
\end{equation}
All transition annuli lie between fixed multiples of $R$; the constants
can depend on the prescribed timelike ADM vector, but not on $R$.
\end{lemma}

\begin{proof}
Use the description $T=v\cdot y$, $y=Ax$, from
Lemma~\ref{n3:reduction:boosted-reference}.  Choose $R_b=c_bR$ so large,
with $c_b$ fixed, that the region $|y|\ge R_b$ is entirely inside the
exact reference part.  Let $\psi$ be a smooth function equal to one on
$(-\infty,0]$, zero on $[1,\infty)$, and taking values in $[0,1]$.
Replace the plane by the graph
\begin{equation}
 \label{n3:reduction:bending-graph}
 T_R(y)=(v\cdot y)
           \psi\left(\frac{\log(|y|/R_b)}{L}\right).
\end{equation}
Its Euclidean gradient satisfies
\[
 |dT_R|_\delta\le |v|\bigl(1+\|\psi'\|_\infty/L\bigr).
\]
Choose a fixed $L$ making the right side strictly smaller than one.
The lapse and spatial coefficients in
Equation~\eqref{n3:reduction:schwarzschild-spacetime} converge to their
Minkowski values, so for all sufficiently large $R$ this graph is
uniformly spacelike.  In fact for positive Schwarzschild mass the lapse
is smaller than one and the spatial conformal factor exceeds one,
which only improves this estimate on the static exterior.

The graph agrees with the original plane before $R_b$ and with a
constant-time slice after $e^LR_b$.  Its first derivative is bounded,
and its second and third derivatives are $O(r^{-1})$ and $O(r^{-2})$.
The induced metric is uniformly comparable with $\delta$; differentiating
its pullback formula gives the first bound in
Equation~\eqref{n3:reduction:bend-geometry}.  The second fundamental form
uses the graph's second derivatives and ambient first derivatives,
divided by a uniformly positive spacelikeness factor, and gives the
other bound.  These estimates remain true in the $x$ coordinates because
$A$ is a fixed invertible linear map.  The constraints vanish on the
entire bent part by the Gauss--Codazzi equations in exact vacuum
Schwarzschild spacetime.

It remains to make the radius precise.  On the bent region keep
$\rho_R=|x|=|A^{-1}y|$.  Beyond the bend, where the slice is already
static, write $y=r\omega$ and put $a(\omega)=|A^{-1}\omega|$.  Choose
$R_c$ a sufficiently large fixed multiple of $R_b e^L$, and a cutoff
$\chi_0$ equal to one before zero and zero after one.  In that static
region set
\[
 \rho_R(r,\omega)
 =r\exp\left[
    \chi_0\left(\frac{\log(r/R_c)}{L_0}\right)\log a(\omega)\right].
\]
For fixed sufficiently large $L_0$ its radial derivative is bounded
below by a positive constant: differentiating in $r$ gives the positive
factor $\rho_R/r$ times
$1+L_0^{-1}\chi_0'\log a$, which is at least $1/2$.
The function and its first two Euclidean derivatives have the bounds
$\rho_R\asymp r$, $|d\rho_R|\le C$, and
$|\partial^2\rho_R|\le C/r$.  Before this interpolation, the same
bounds hold for $|A^{-1}y|$, and uniform metric comparability gives the
gradient lower bound.  Combining these facts with
Equation~\eqref{n3:reduction:bend-geometry} proves
Equation~\eqref{n3:reduction:radius-geometry}.  After the interpolation,
$\rho_R=r$ is precisely the static isotropic radius.  All radius ratios
used in the construction are fixed, independently of $R$.
\end{proof}

\begin{lemma}[A radial DEC repair of vanishing mass cost]
\label{n3:reduction:radial-repair}
Consider the data obtained from
Proposition~\ref{n3:reduction:annular-replacement} and
Lemma~\ref{n3:reduction:vacuum-bend}.  There is a smooth positive conformal
factor $u_R$, constant on the unchanged inner region and tending to one
at infinity, such that $(u_R^4g_R,u_R^2K_R)$ satisfy DEC everywhere.
Moreover
\begin{equation}
 \label{n3:reduction:repair-size}
 \|u_R-1\|_\infty=o(R^{-1}),\qquad
 |du_R|=o(R^{-2})\quad\hbox{on the transition region}.
\end{equation}
The final end is exactly Schwarzschild with $K=0$ and mass
$m+o(1)$.  The boundary expansion retains its original weak sign.
\end{lemma}

\begin{proof}
All geometry estimates below are uniform in $R$.  Enlarge the fixed
transition region slightly so that, in terms of $t=\rho_R/R$, it lies
inside an interval $(a,b)$ with $0<a<b<\infty$.  Arrange that $t\ge b$
is in the static spherical region, and that the support of the possible
deficit lies in a fixed smaller interval $I\Subset(a,b)$.  Increasing
$b$ if needed makes these assertions follow from
Lemma~\ref{n3:reduction:vacuum-bend}.
We take $a=1$: the compact supports in the annular correction leave
a fixed neighborhood of $t=1$ free of deficit.  The factor constructed
below will consequently be constant throughout the original region
$r\le R$, and is extended by that constant over the compact core of $N$.

Let $\eta_R\to0$ be as in
Equation~\eqref{n3:reduction:annular-deficit}, replacing it by a positive
majorant if necessary.  We will choose $\varepsilon_R=C_0\eta_R$ and
construct
\[
 u_R=1+\frac{\varepsilon_R}{R}F_R(\rho_R/R),
\]
where $F_R$ and the relevant derivatives are bounded independently of
$R$.  Write $z_R=-F_R'$.  On $[a,b+1]$, initially define $z_R$ by
\begin{equation}
 \label{n3:reduction:repair-ode}
 z_R'=C_1z_R+\omega_1,
 \qquad z_R(a)=0,
\end{equation}
where $\omega_1$ is smooth, nonnegative, zero near $a$, and at least one
on $I$.  Take $C_1$ sufficiently large, depending only on the constants
in Equations~\eqref{n3:reduction:bend-geometry} and
\eqref{n3:reduction:radius-geometry}.  The solution is nonnegative and
has uniform bounds on this fixed interval.  For a function with
$F_R'=-z_R$ the chain rule gives
\begin{align*}
 -\Delta_{g_R}u_R-|K_R|_{g_R}|du_R|_{g_R}
 &=\frac{\varepsilon_R}{R^3}
   \left[z_R'|d\rho_R|_{g_R}^2
       +Rz_R\Delta_{g_R}\rho_R
       -Rz_R|K_R|_{g_R}|d\rho_R|_{g_R}\right]\\
 &\ge\frac{\varepsilon_R}{R^3}
       (c_1z_R'-C_2z_R).
\end{align*}
Choose $C_1$ so that this is nonnegative and is at least
$c_1\varepsilon_RR^{-3}$ on $I$.

The preceding differential inequality need not be imposed on the
spherical tail.  There $K_R=0$ and the Schwarzschild radial Laplace flux
coefficient is $(r+m/2)^2$.  Put $c_R=m/(2R)$ and, on the still spherical
interval near $b$, define
\[
 A_R(t)=(t+c_R)^2z_R(t).
\]
It has positive derivative wherever $z_R>0$.  Continue it to a positive
constant on $[b+1,\infty)$ while keeping $A_R'\ge0$.  For an explicit
smooth continuation, continue the solution of
Equation~\eqref{n3:reduction:repair-ode} to $b+1$, multiply the derivative
of $(t+c_R)^2z_R(t)$ by a nonnegative cutoff equal to one near $b$ and
zero near $b+1$, and integrate with the initial value $A_R(b)$.
Define $z_R=A_R(t)/(t+c_R)^2$ on this tail, and set it to zero for
$t\le a$.  All joins are smooth, the functions are nonnegative, and
their bounds on finite intervals are uniform in $R$.  The constant
tail value, denoted $A_R^\infty$, is bounded independently of $R$.

Now define
\[
 F_R(t)=\int_t^\infty z_R(s)\,\dd s.
\]
It is nonnegative, uniformly bounded, and constant before $a$.
The preceding construction gives a superharmonic $u_R$ on the spherical
tail, because its signed radial flux is
$-\varepsilon_RA_R(t)$ and is nonincreasing.  It also gives
Equation~\eqref{n3:reduction:repair-size}.  For large $R$, $1\le u_R\le2$.
Increasing $C_0$ ensures
\[
 -\Delta_{g_R}u_R-|K_R||du_R|
 \ge \frac{u_R}{8}\,
       \bigl[16\pi(|J_R|_{g_R}-\mu_R)\bigr]_+.
\]
On the original inner region the conformal factor is constant and the
original data satisfy DEC.  Elsewhere this inequality and
Lemma~\ref{n3:reduction:conformal-formulas} prove DEC for the repaired data.
Using $u_R$, rather than estimating $\log u_R$ separately, incorporates
the gradient-square term exactly.

On the final end,
\[
 u_R=1+\frac{\varepsilon_RA_R^\infty}{r+m/2}.
\]
Writing $a_R=\varepsilon_RA_R^\infty=o(1)$, one obtains the exact identity
\begin{equation}
 \label{n3:reduction:repaired-schwarzschild}
 u_R^4\left(1+\frac{m}{2r}\right)^4\delta
   =\left(1+\frac{m/2+a_R}{r}\right)^4\delta.
\end{equation}
Thus the final mass is $m+2a_R=m+o(1)$ and the momentum is zero.
Multiplication of $K_R$ by $u_R^2$ preserves its compact support.
On the inner boundary, where $u_R$ is constant, the expansion is simply
multiplied by $u_R^{-2}$, so its weak sign is preserved.
\end{proof}

\subsection{The numerical reduction}

The preceding construction has replaced the end while preserving DEC
and the boundary expansion sign.  A final strict conformal change puts
the data in the class used for the elliptic deformation.

\begin{proposition}[Reduction to strict data with a controlled rest end]
\label{n3:reduction:rest-reduction}
Let $(N,g,K)$ satisfy the exterior hypotheses of
Theorem~\ref{n3:intro:exterior-inequality}, including $\theta_+\le0$ on its
smooth compact boundary, and let $m=(E^2-|P|^2)^{1/2}>0$.
There is a sequence of smooth exterior data $(g_j,K_j)$ on $N$ with the
following properties:
\begin{enumerate}[label=\textup{(\roman*)}]
\item $\mu_j>|J_j|_{g_j}$ everywhere and $\theta_{+,j}<0$ on $B$.
      The constraint densities are integrable.  On the end their DEC
      margin is bounded below by $c_jr^{-3-\delta}$ for some $c_j>0$
      and one fixed $0<\delta<1$.
\item $K_j$ is compactly supported, and in a rest chart
      $g_j-\delta=O_k(r^{-1})$ and $R_{g_j}=O_k(r^{-3-\delta})$
      for every $k$.  In particular
      $\Ric_{g_j}=O(r^{-3})$ there.  Their ADM momentum is zero.
\item $(g_j,K_j)\to(g,K)$ smoothly on every fixed compact subset up to
      the boundary, their ADM energies satisfy $E_j\to m$, and
      $a_{g_j}(B)\to a_g(B)$.
\end{enumerate}
Consequently it suffices to prove the energy/enclosing-area inequality
for the class in \textup{(i)}--\textup{(ii)} in order to prove
Theorem~\ref{n3:intro:exterior-inequality} in its stated class.
\end{proposition}

\begin{proof}
Apply Proposition~\ref{n3:reduction:annular-replacement},
Lemma~\ref{n3:reduction:vacuum-bend}, and
Lemma~\ref{n3:reduction:radial-repair} along any sequence $R_j\to\infty$.
The resulting DEC exteriors have compactly supported second fundamental
form and an exact rest Schwarzschild end with energy tending to $m$.
They agree with the original data on expanding compact subsets except
for a constant conformal factor tending to one.  Hence they converge
smoothly on each fixed compact subset.  Their boundary expansion is
weakly negative.

Apply the last part of Lemma~\ref{n3:reduction:strict-direction} separately
to each of these fixed exteriors, using the same chosen
$0<\delta<1$ and a compactly supported drift coefficient.  Choose its
positive parameter $s$ in Equation~\eqref{n3:reduction:linear-family}
small enough that the energy change, $\|\log(1+s\phi)\|_\infty$, and
the supremum of its first derivative multiplied by the fixed original
coordinate radius are at most $1/j$.  Also require the change in the
first $j$ smooth norms on the first $j$ members of a compact exhaustion
to be at most $1/j$.
All these quantities are finite for the fixed $j$-th exterior, so these
choices are possible.  No estimate uniform in the receding bending
radius is required for this final strictification.

The resulting data satisfy \textup{(i)}. The strict conformal factor has
symbol decay of all orders and does not enlarge the support of $K$.
On the far Schwarzschild end the drift vanishes, so
$-\Delta\phi=w=r^{-3-\delta}$. Since the rest Schwarzschild metric
has zero scalar curvature and $K=0$ there, the exact $u$-formula gives
\[
 R_{g_j}=8s(1+s\phi)^{-5}r^{-3-\delta}.
\]
This proves the scalar-curvature bound in \textup{(ii)}.
The Ricci bound follows from
the two-derivative metric decay. The energy change tends to zero and the
momentum change is zero. The prescribed smallness gives smooth local
convergence, proving \textup{(ii)} and the local convergence and charge
assertions in \textup{(iii)}.

The metrics before strictification satisfy the uniform comparisons and
scaled first-derivative bounds of
Lemma~\ref{n3:reduction:area-compactness}, in the fixed original end chart:
the reference plane, the bend, and the radius interpolation involve only
fixed linear maps and fixed radius ratios.  The radial repair preserves
these bounds.  The final strictification parameters were chosen to
preserve them as well.  That lemma therefore gives the convergence of
enclosing infima.

If the reduced energy inequality is available, applying it to these
data gives
\[
 E_j\ge\sqrt{a_{g_j}(B)/(16\pi)}.
\]
Passing to the limit proves
$m\ge\sqrt{a_g(B)/(16\pi)}$.  Only the energies and enclosing infima
enter this passage to the original data.
\end{proof}
